\documentclass{article}
\usepackage{amsmath}    % For advanced mathematical typesetting
\usepackage{amssymb}    % For additional mathematical symbols
\usepackage{enumitem}   % For customizable lists (e.g., itemize with bold labels)
\usepackage[numbers,sort&compress]{natbib} % For citation management (using natbib for flexibility, numbers style)
\usepackage{url}        % To handle URLs in bibliography
\usepackage{hyperref}   % For clickable links (optional but good practice)
\usepackage{booktabs}   % For better table rules
\usepackage{tabularx}   % For tables with adjustable width columns
\usepackage{algorithm}  % For algorithm environments
\usepackage{algpseudocode} % For pseudocode within algorithm
\usepackage[margin=1in]{geometry} % Added geometry package for margin control
\usepackage{authblk}    % Added package for author affiliations
\usepackage{tikz}       % Added for drawing diagrams
\usepackage{graphicx}   % Added for including images
\usetikzlibrary{shapes.geometric, arrows.meta, positioning, backgrounds, fit, shapes.symbols} % Added TikZ libraries

% Simple Keywords command definition
\newcommand{\keywords}[1]{\textbf{Keywords: } #1}


\hypersetup{
    colorlinks=true,
    linkcolor=blue,
    filecolor=magenta,      
    urlcolor=cyan,
    citecolor=green, % Make citations green for visibility
}

\bibliographystyle{ieeetr} % IEEE transaction style bibliography

\title{Self-Defending 6G Networks Through AI-Driven Adaptive Decoy Generation at the Edge}

% Authors and Affiliations using authblk package
\author[1]{Ranil Mukesh MJ}
\author[2]{Pandiya Rajan G}
\author[1]{Prabhu Gopal}
\author[4]{Malathy Sathyamoorthy\thanks{Corresponding author: \texttt{ksmalathy@gmail.com}}}
\author[4]{Aniket S. Nagane}
\author[5]{Rakesh Keshava}
\affil[1]{\href{https://phobosq.com}{PhobosQ Private Limited , Coimbatore , India}}
\affil[2]{Department of Computer Science and Engineering (Artificial Intelligence and Machine Learning), KPR Institute of Engineering and Technology, Coimbatore, India}
\affil[3]{Department of Information Technology, KPR Institute of Engineering and Technology, Coimbatore, India}
\affil[4]{Symbiosis Institute of Computer Studies and Research (SICSR), Symbiosis International (Deemed University), Pune, India}
\affil[5]{Security Engineering, California, USA 94555}

% \date{\today} % Or specify a date

\begin{document}

\maketitle % Display title/author if defined

\begin{abstract}
As 6G networks introduce unprecedented connectivity and a vastly expanded attack surface, traditional security paradigms prove insufficient. This paper presents a novel Adaptive Decoy Generation framework that provides proactive, intelligent, and self-evolving defense at the network edge. The framework integrates three core AI components: a Conditional Generative Adversarial Network (cGAN) to synthesize high-fidelity, context-aware decoy entities; a Reinforcement Learning (RL) agent guided by Proximal Policy Optimization (PPO) to dynamically manage decoy deployment and resource allocation; and a unique adversarial feedback loop that enables the system to learn from attacker interactions and continuously adapt its strategies. Deployed within isolated network slices and synchronized via federated learning to ensure privacy and scalability, the framework is optimized for 6G's distributed architecture.Simulation results, using the CIC-IoT-2023 dataset, demonstrate the framework's exceptional performance, achieving a 98.2\% detection rate with a near-zero false positive rate and an average detection latency of approximately 15 ms at the edge. Crucially, the system exhibits a positive Adaptability Index (+5.7\%), confirming its ability to counter sophisticated, learning-based attackers a significant advantage over static defense mechanisms. This work establishes a comprehensive blueprint for a resilient, AI-driven security architecture, paving the way for self-defending 6G ecosystems.
\end{abstract}

\keywords{6G Networks, Large Language Models (LLMs), Proximal Policy Optimization (PPO), Network Security, Anomaly Detection, Generative AI, Generative Adversarial Networks (GANs), Conditional GANs, Decoy Generation, Proactive Security, Edge Computing, Reinforcement Learning, Federated Learning} % Expanded LLMs and PPO

\section{Introduction}

Sixth-generation (6G) wireless networks enable hyper-connected ecosystems characterized by unprecedented data rates, ultra-low latency, massive device density, and the native integration of artificial intelligence (AI) \cite{Imoize20216GESX, Li2025EvolvingTAL}. These advancements support transformative applications across domains such as autonomous systems, immersive extended reality (XR), holographic communication, and the Internet of Everything (IoE). However, the complexity, heterogeneity, and expanded attack surface of 6G architectures defined by distributed edge computing, software-defined networking (SDN), network slicing, and AI-driven operations introduce significant security and privacy challenges \cite{Nguyen2021SecurityAPN, Hussein2022SecurityRAAA, Li2024TheSAR, Abuadbba2024From5TAM, Salas2024SecurityTAAN, Janay2023SecureDTAO}. Traditional security mechanisms, which are reactive and signature-based, do not address the scale, speed, and sophistication of threats in the 6G era \cite{Rahman2021NetworkAPAQ, Siriwardhana2021AIA6W}.

Artificial intelligence, particularly machine learning (ML) and deep learning (DL), strengthens the security of 6G networks \cite{Haider2020ArtificialIAT, Siriwardhana2021AIA6W, Li2025EvolvingTAL}. AI enhances network automation, resource optimization, and security functions such as intrusion detection and anomaly analysis \cite{Cifuentes2024AnalysisOTE, Gkonis2024LeveragingNDAF, Rodrguez2024ASSQ}. Among advanced AI techniques, Generative Adversarial Networks (GANs) model complex data distributions, generate synthetic data, and identify deviations from normal behavior \cite{Vu2024ApplicationsOGC, Ayanoglu2022MachineLIA}. These capabilities enable advanced anomaly detection within the dynamic, high-dimensional data landscape of 6G \cite{Ferrag2023GenerativeAFB, Arifin2024ASOF}.

The Adaptive Decoy Generation method utilizes Conditional GANs (cGANs) to proactively detect and analyze threats within 6G networks. The system actively generates and deploys realistic, context-aware decoy entities such as devices, services, or traffic patterns at the network edge. Interactions with these decoys are immediately flagged as malicious, enabling early detection and intelligence gathering. An adversarial learning feedback loop uses data from attacker interactions to continuously refine the generative model, improving decoy realism and adaptability against evolving threats. Reinforcement Learning (RL) supports dynamic decoy management, and key 6G features such as edge computing and network slicing enable efficient and isolated deployment.

The methodology encompasses cGAN-based decoy generation, edge optimization techniques, RL-driven dynamic deployment strategies, real-time anomaly detection mechanisms, an adversarial feedback mechanism, and federated learning for model synchronization. The technical blueprint is supported by mathematical formulations and implementation considerations, establishing a practical and scalable solution for enhancing the security posture of next-generation wireless networks. Subsequent sections review relevant literature, present the methodology, discuss evaluation metrics, and outline the impact of this research.

\section{Literature Review}

The security landscape of future wireless networks, particularly 6G, is rapidly evolving. The unique characteristics of 6G, including terahertz communications, an AI-native architecture, and massive interconnectivity, necessitate the development of novel security paradigms \cite{Nguyen2021SecurityAPN, Hussein2022SecurityRAAA, Li2024TheSAR, Abuadbba2024From5TAM, Salas2024SecurityTAAN}. The research community broadly recognizes that conventional security measures are insufficient, highlighting the need for intelligent, adaptive, and proactive defense mechanisms \cite{Siriwardhana2021AIA6W, Janay2023SecureDTAO, Rahman2021NetworkAPAQ}.

Artificial intelligence and machine learning have become central to enhancing network security within 5G and emerging 6G systems \cite{Haider2020ArtificialIAT, Li2025EvolvingTAL}. Recent studies have explored the application of AI for intrusion detection, threat prediction, and automated security orchestration \cite{Cifuentes2024AnalysisOTE, Gkonis2024LeveragingNDAF, Rodrguez2024ASSQ, Yang2023TowardsACK}. Anomaly detection remains a critical area in which AI demonstrates significant capability. Techniques such as deep neural networks (DNNs), convolutional neural networks (CNNs), and recurrent neural networks (RNNs) have been applied to identify anomalous traffic patterns in network data \cite{Saeed2023AnomalyDIY, Rzym2024DynamicTAAB, Solanki2024EnhancedADAL, Cao2024TheDOH, Paolini2023RealTimeCBP, Nomikos2025ADTAI}. However, these methods often rely on labeled datasets, which may be scarce or outdated, and can be limited when addressing zero-day attacks or sophisticated adversarial manipulations.

Generative AI, particularly GANs, has attracted considerable attention for its applications in networking and cybersecurity \cite{Vu2024ApplicationsOGC, Ayanoglu2022MachineLIA, Arifin2024ASOF}. GANs learn the underlying distribution of complex, high-dimensional data, making them suitable for generating synthetic data, augmenting datasets, and modeling communication channels \cite{Wen2024GenerativeAFU, Tao2023WirelessNDZ, Zhang2023GenerativeAFS, Muhammad2024IntegratingGAAE}. In security, GANs have been used to generate adversarial examples for evaluating detector robustness, to enhance intrusion detection systems, and to detect anomalies \cite{Park2023AnEAD, Rao2024AIDAG, Gonzalez2020OnTUJ, Ferrag2023GenerativeAFB}. For example, Park et al. \cite{Park2023AnEAD} developed a GAN-based Network Intrusion Detection System (NIDS) that demonstrated improved performance. Rao et al. \cite{Rao2024AIDAG} employed a hybrid CNN-GAN model for anomaly detection in network traffic. Gonz\'{a}lez et al. \cite{Gonzalez2020OnTUJ} investigated generative models for anomaly detection in multivariate network time-series data.

The use of generative models for proactive defense strategies, such as decoy or honeypot generation, is an emerging research area. While traditional honeypots exist, AI-driven decoys provide greater realism and adaptability. Ferrag et al. \cite{Ferrag2023GenerativeAFB} discussed generative AI for cyber threat hunting in 6G IoT networks, aligning with a proactive security stance. Adaptive security, in which defense mechanisms evolve based on observed threats, is recognized as essential for dynamic environments such as 6G \cite{Shen2021AdaptiveADAC, Vemuri2024AdaptiveGAM}.

The integration of AI with 6G architectural features, including edge computing and federated learning, is critical for practical deployment. Edge computing enables low-latency processing for real-time threat detection \cite{Ferrag2023EdgeLFAD}, while federated learning supports collaborative model training without sharing raw data, thus preserving privacy \cite{Das2022FGANFGI}. Das \cite{Das2022FGANFGI} introduced Federated GANs (FGAN) for anomaly detection, demonstrating the potential of distributed generative models. Reinforcement learning has also been applied for dynamic network management and security optimization \cite{Shen2021AdaptiveADAC}.

Although research has addressed GANs for anomaly detection \cite{Park2023AnEAD, Rao2024AIDAG, Gonzalez2020OnTUJ}, federated learning for security \cite{Das2022FGANFGI}, and the broader role of AI in 6G security \cite{Siriwardhana2021AIA6W, Ferrag2023GenerativeAFB}, a gap remains in integrating these concepts into a cohesive, proactive security framework tailored to the challenges of 6G. Existing approaches often focus on passive detection or lack mechanisms for continuous adaptation based on attacker interactions. The present methodology addresses this gap by combining:
\begin{itemize}
    \item Conditional GANs for generating context-aware, highly realistic decoys.
    \item Strategic deployment at the 6G edge, optimized for resource constraints.
    \item Reinforcement learning for dynamic, intelligent management of decoy resources.
    \item An adversarial feedback loop that uses attacker interaction data to refine decoy generation, enhancing long-term effectiveness and adaptability.
    \item Integration with core 6G features such as network slicing and federated learning for isolation and scalable model synchronization.
\end{itemize}
This integrated approach provides a robust, adaptive, and proactive security solution for 6G networks. It advances beyond conventional detection paradigms to enable active threat engagement and intelligence gathering. Concerns regarding the trustworthiness and privacy implications of generative AI \cite{Li2024TrustworthyACO, Feretzakis2024PrivacyPreservingTIAG} are acknowledged, and the use of federated learning and network slicing within this framework contributes to mitigating these issues.

\section{Methodology}

The Adaptive Decoy Generation method establishes a security framework for 6G networks by leveraging Conditional Generative Adversarial Networks (cGANs) to synthesize realistic decoy entities such as devices, services, or traffic patterns that emulate legitimate network components \cite{Ayanoglu2022MachineLIA, Park2023AnEAD}. These decoys are strategically deployed at edge nodes to attract and detect attackers, with any interaction flagged as anomalous. Unlike conventional passive anomaly detection techniques \cite{Saeed2023AnomalyDIY}, this method engages threats directly, enabling early detection and intelligence gathering \cite{Ferrag2023GenerativeAFB}. An adversarial learning feedback loop refines the generative model using attacker interactions, enhancing decoy realism and adaptability against evolving threats \cite{Vemuri2024AdaptiveGAM}.

This approach is designed for 6G's characteristics: ultra-high data rates, low latency, distributed edge computing, and an AI-native architecture \cite{Li2025EvolvingTAL, Imoize20216GESX}. A reinforcement learning (RL)-driven decoy management system dynamically optimizes deployment based on real-time network conditions and threat levels \cite{Shen2021AdaptiveADAC}. The methodology ensures scalability, efficiency, and robustness, providing a practical solution for securing next-generation networks \cite{Rodrguez2024ASSQ}. The following sections detail the technical blueprint, including decoy generation, deployment, anomaly detection, and synchronization, supported by mathematical formulations and implementation considerations. Figure~\ref{fig:system_architecture} illustrates the overall system architecture.

% <<< Figure 1 image replaced >>>
\begin{figure}[htbp]
    \centering
    \includegraphics[width=0.8\textwidth]{fig1NEW.png}
    \caption{Simplified system architecture: Central server coordinates edge nodes. Edge nodes deploy decoys (in a decoy slice) to attract attackers and real services (in a legitimate slice) for users. Attacker interactions with decoys are fed back to the central server for continuous improvement.}
    \label{fig:system_architecture}
\end{figure}


\subsection{Decoy Generation Using Conditional GANs}
\subsubsection{Model Selection: Conditional GANs (cGANs)}
To generate diverse, context-specific decoys (e.g., IoT devices, vehicular nodes), Conditional Generative Adversarial Networks (cGANs) are employed \cite{Ayanoglu2022MachineLIA}. In cGANs, the generator $G$ and discriminator $D$ are conditioned on auxiliary information $y$ (e.g., device type, service category), enabling tailored decoy synthesis. The training objective is formulated in Equation~\eqref{eq:cgan_objective}:
\begin{equation}
\label{eq:cgan_objective}
\min_G \max_D \mathbb{E}_{x \sim p_{\mathrm{data}}, y \sim p_y} [\log D(x, y)] + \mathbb{E}_{z \sim p_z, y \sim p_y} [\log (1 - D(G(z, y), y))]
\end{equation}
Here, $G(z, y)$ produces synthetic decoys conditioned on $y$, while $D(x, y)$ assesses their realism within that specific context. Equation~\eqref{eq:cgan_objective} ensures that generated decoys are both realistic and diverse, which is critical for deceiving attackers within 6G's heterogeneous environment \cite{Vu2024ApplicationsOGC}.

\subsubsection{Training Process with Enhanced Stability}
The cGAN is trained using a dataset comprising legitimate 6G traffic, represented as $\mathcal{D}_{\mathrm{train}} = \{(x_1, y_1), \dots, (x_n, y_n)\}$. During training, the generator learns the conditional data distribution $p_{\mathrm{data}}(x|y)$. The Wasserstein GAN with Gradient Penalty (WGAN-GP) objective \cite{Arifin2024ASOF} is used to enhance training stability. The procedure involves iterative updates to the discriminator and generator, as detailed in Algorithm~\ref{alg:gan_training}.

\begin{algorithm}
\caption{cGAN Training with WGAN-GP}
\label{alg:gan_training}
\begin{algorithmic}[1]
\Require Legitimate dataset \(\mathcal{D}_{\mathrm{train}}\), learning rate \(\alpha\), batch size \(m\), number of critic iterations \(n_{\mathrm{critic}}\), gradient penalty coefficient \(\lambda\)
\State Initialize generator \(G_\theta\) and discriminator \(D_\omega\) parameters \(\theta, \omega\)
\While{not converged}
    \For{\(t=1\) to \(n_{\mathrm{critic}}\)}
        \State Sample minibatch of real data \(\{(x^{(i)}, y^{(i)})\}_{i=1}^m \sim p_{\mathrm{data}}\)
        \State Sample minibatch of latent vectors \(\{z^{(i)}\}\}_{i=1}^m \sim p_z\) and associated conditions \(\{y^{(i)}\}\}_{i=1}^m \sim p_y\)
        \State Generate fake samples \(\tilde{x}^{(i)} = G_\theta(z^{(i)}, y^{(i)})\)
        \State Sample \(\epsilon \sim U[0, 1]\)
        \State Calculate interpolated sample \(\hat{x}^{(i)} = \epsilon x^{(i)} + (1-\epsilon) \tilde{x}^{(i)}\)
        \State Calculate gradient penalty \( P = (\|\nabla_{\hat{x}} D_\omega(\hat{x}^{(i)}, y^{(i)})\|_2 - 1)^2 \)
        \State Calculate discriminator loss \( L_D = \frac{1}{m}\sum [D_\omega(\tilde{x}^{(i)}, y^{(i)}) - D_\omega(x^{(i)}, y^{(i)})] + \frac{\lambda}{m}\sum P \)
        \State Update discriminator parameters \( \omega \leftarrow \omega - \alpha \cdot \nabla_\omega L_D \) (using Adam or RMSProp)
    \EndFor
    \State Sample minibatch of latent vectors \(\{z^{(i)}\}\}_{i=1}^m \sim p_z\) and conditions \(\{y^{(i)}\}\}_{i=1}^m \sim p_y\)
    \State Calculate generator loss \( L_G = -\frac{1}{m}\sum D_\omega(G_\theta(z^{(i)}, y^{(i)}), y^{(i)}) \)
    \State Update generator parameters \( \theta \leftarrow \theta - \alpha \cdot \nabla_\theta L_G \) (using Adam or RMSProp)
\EndWhile
\end{algorithmic}
\end{algorithm}

This WGAN-GP approach aids in mitigating issues such as mode collapse and ensures the generation of higher quality decoys \cite{Park2023AnEAD}.

\subsubsection{Edge-Optimized Generative Models}
Considering 6G's significant reliance on edge computing \cite{Ferrag2023EdgeLFAD}, the cGAN model is optimized for resource-constrained edge nodes using techniques discussed in contexts such as efficient AI deployment \cite{Zhou2024GenerativeAAAH}:
\begin{itemize}[label=\textbf{\textbullet}]
    \item \textbf{Knowledge Distillation}: A lightweight "student" generator is trained to mimic the behavior of the full cGAN, thereby reducing computational overhead.
    \item \textbf{Quantization-Aware Training (QAT)}: Model weights are quantized (e.g., converted to 8-bit integers) during the training process, preserving performance when deployed on edge devices.
    \item \textbf{Sparsity-Inducing Pruning}: Non-critical neural network connections are removed (pruned), resulting in a more efficient model suitable for deployment.
\end{itemize}
These optimizations facilitate real-time decoy generation directly at the network edge, aligning effectively with 6G's distributed computing paradigm \cite{Gkonis2024LeveragingNDAF}.

\subsection{Decoy Deployment and Dynamic Management}
\subsubsection{Distributed Deployment at Edge Nodes}
Decoys are generated and deployed locally on edge nodes, utilizing 6G's edge computing infrastructure to minimize latency and reduce central processing overhead \cite{Ferrag2023EdgeLFAD}. This distributed strategy scales effectively with 6G's expected massive connectivity, ensuring comprehensive coverage across diverse network segments \cite{Imoize20216GESX}.

\subsubsection{Isolation via Network Slicing}
Decoys are isolated within dedicated virtualized network slices, a fundamental feature of 6G \cite{Nguyen2021SecurityAPN}, to prevent interference with legitimate network operations and contain potential security compromises. Each decoy or group operates within its own slice, confining interactions predominantly to attackers and thereby mitigating the risk of false positives affecting legitimate users \cite{Abuadbba2024From5TAM}.

\subsubsection{Reinforcement Learning for Dynamic Management}
A reinforcement learning (RL) agent is employed to dynamically optimize decoy deployment strategies. The agent interacts with the environment (defined by the current 6G network state and threat landscape) by following a policy \(\pi(a|s)\) designed to maximize cumulative rewards. The state \(s\), action \(a\), and reward \(R\) components are defined as follows:
\begin{itemize}[label=\textbf{\textbullet}]
    \item \textbf{State} \( s \): Encompasses network telemetry data (e.g., traffic load, active service types), historical attack patterns, decoy interaction rates, and current resource availability at edge nodes.
    \item \textbf{Action} \( a \): Consists of decisions such as activating or deactivating specific decoy types, adjusting the density of deployed decoys, or modifying decoy parameters (e.g., simulated operating system, open ports).
    \item \textbf{Reward} \( R \): Represents a balance between successful threat detection (measured by unique interactions flagged as malicious) and resource efficiency (considering computational and network overhead).
\end{itemize}
The primary objective is to determine an optimal policy \(\pi^*\) that maximizes the expected discounted cumulative reward as shown in Equation~\eqref{eq:rl_objective}:
\begin{equation}
\label{eq:rl_objective}
\max_\pi \mathbb{E} \left[ \sum_{t=0}^{\infty} \gamma^t R(s_t, a_t) \right]
\end{equation}
where \(\gamma\) is the discount factor controlling the importance of future rewards. We propose utilizing Proximal Policy Optimization (PPO) \cite{Yang2023TowardsACK}, known for its stability and efficiency as an actor-critic algorithm, making it well-suited for this context. A conceptual sketch of the PPO process is provided in Algorithm~\ref{alg:ppo_rl}.

\begin{algorithm}
\caption{Dynamic Decoy Management using PPO}
\label{alg:ppo_rl}
\begin{algorithmic}[1]
\Require Initialize policy \(\pi_\theta\) and value function \(V_\phi\) parameters \(\theta, \phi\)
\For{each iteration}
    \State Collect set of trajectories \(\mathcal{T}_k = \{(s_t, a_t, R_t)\}\) by executing policy \(\pi_\theta\) within the environment (edge network).
    \State Compute rewards-to-go \(\hat{R}_t\) and advantage estimates \(\hat{A}_t\) based on the current value function \(V_\phi\).
    \For{each epoch}
        \State Optimize the surrogate objective for the policy network using stochastic gradient ascent:
        \State \( \theta \leftarrow \theta + \alpha_\theta \nabla_\theta \mathcal{L}^{\text{CLIP}}(\theta) \)
        \State Optimize the value function loss using gradient descent:
        \State \( \phi \leftarrow \phi - \alpha_\phi \nabla_\phi \mathcal{L}^{\text{VF}}(\phi) \)
    \EndFor
\EndFor
\end{algorithmic}
\end{algorithm}

This RL-based approach ensures adaptive decoy management capable of responding effectively to the fluctuating network conditions characteristic of 6G environments \cite{Shen2021AdaptiveADAC}.

\subsection{Anomaly Detection and Adversarial Feedback Loop}
\subsubsection{Real-Time Interaction Monitoring}
Any interaction directed towards a deployed decoy (including, but not limited to, connection attempts, port scanning, data probes, and authentication attempts) is considered inherently suspicious and consequently flagged as an anomaly \cite{Ferrag2023GenerativeAFB}. This detection mechanism operates directly at the edge node hosting the decoy, leveraging 6G's low-latency communication capabilities to enable near real-time alerting \cite{Rzym2024DynamicTAAB}. All flagged events are logged with comprehensive metadata (e.g., timestamp, source IP address or identifier, type of interaction, targeted decoy profile, and any data exchanged) to facilitate subsequent analysis and feedback integration.

\subsubsection{Adversarial Feedback for Model Refinement}
A novel adversarial feedback loop, illustrated in Figure~\ref{fig:feedback_loop}, is implemented to continuously enhance the effectiveness of the generative model by utilizing data gathered from actual attacker interactions:

% <<< Figure 2 image replaced >>>
\begin{figure}[htbp]
    \centering
    \includegraphics[width=0.6\textwidth]{fig2NEW.png}
    \caption{Simplified adversarial feedback loop: Attacker probes decoy, interactions are logged and aggregated as attack data, which is used to update the generator. Improved decoys are then redeployed, creating a continuous adaptation cycle.}
    \label{fig:feedback_loop}
\end{figure}


\begin{itemize}[label=\textbf{\textbullet}]
    \item \textbf{Interaction Dataset Collection}: Successfully logged interactions with decoys constitute a dataset \(\mathcal{D}_{\mathrm{attack}} = \{x_{\mathrm{attack}}^{(1)}, x_{\mathrm{attack}}^{(2)}, \dots\}\). This dataset captures the specific characteristics of traffic or behaviors exhibited by real attackers when targeting the deployed decoys.
    \item \textbf{Generator Fine-Tuning}: The generator \( G \) undergoes periodic fine-tuning, incorporating insights derived not only from legitimate data but also from the attacker interaction dataset \(\mathcal{D}_{\mathrm{attack}}\). This process can involve introducing an additional loss term that encourages the generator to produce decoys possessing features that proved attractive to attackers. This might be formulated as minimizing the distance between generated decoys and observed attack patterns, conditioned on relevant context \( y \), as expressed in Equation~\eqref{eq:attack_loss}:
    \begin{equation}
    \label{eq:attack_loss}
    \mathcal{L}_{\mathrm{attack}} = \mathbb{E}_{x_{\mathrm{attack}} \sim \mathcal{D}_{\mathrm{attack}}} \left[ \min_{z} \| G(z, y_{\mathrm{context}}) - x_{\mathrm{attack}} \|^2 \right]
    \end{equation}
    where \( y_{\mathrm{context}} \) represents the condition pertinent to the specific attack context. Equation~\eqref{eq:attack_loss} assists the generator in adapting to attacker preferences and evolving reconnaissance techniques.
    The minimization over the latent vector \( z \) required in \(\mathcal{L}_{\mathrm{attack}}\) constitutes a GAN inversion task. For each observed attack sample \( x_{\mathrm{attack}} \), the goal is to identify the latent code \( z^* \) that best reconstructs the sample via the generator: \( z^* = \arg\min_{z} \| G(z, y_{\mathrm{context}}) - x_{\mathrm{attack}} \|^2 \). This optimization can be achieved using iterative gradient-based methods (e.g., executing approximately 50 steps of Adam optimization directly on \( z \), potentially clamping values to a valid range like \([-1, 1]\)) or by training a separate encoder network designed to approximate the inverse mapping \cite{Feng2022NearPG}.
    \item \textbf{Discriminator Update}: The discriminator \( D \) can also be updated using samples from \(\mathcal{D}_{\mathrm{attack}}\) (these might be labeled as 'fake' or assigned a distinct 'attack' class). This improves its ability not only to distinguish generated decoys from legitimate data but potentially also to recognize attacker-specific patterns, thereby providing further guidance to the generator.
\end{itemize}
This adaptive feedback loop ensures that the decoys evolve dynamically in response to observed attack strategies, thereby maintaining their effectiveness against adversaries who might otherwise learn to identify and circumvent static decoys over time \cite{Vemuri2024AdaptiveGAM}.

\subsubsection{Alert System and Threat Intelligence}
Anomaly reports, containing concise summaries of detected interactions (e.g., potential attack signatures, source information, targeted decoy type), are transmitted to a central security management system or a Security Operations Center (SOC). This communication process leverages 6G's Ultra-Reliable Low-Latency Communication (URLLC) capabilities where necessary, particularly for critical alerts \cite{Nguyen2021SecurityAPN}. The aggregated data provides a comprehensive, global view of threat activity across the network, enabling broader threat intelligence analysis and facilitating coordinated response actions \cite{Rodrguez2024ASSQ}.

\subsection{Communication and Synchronization}
\subsubsection{Federated Learning for Model Synchronization}
Federated learning (FL) is employed to synchronize the cGAN models across distributed edge nodes without necessitating the sharing of raw, potentially sensitive local data \cite{Das2022FGANFGI, Ferrag2023EdgeLFAD}. The standard Federated Averaging (FedAvg) process is adapted for this purpose, as illustrated in Algorithm~\ref{alg:fl_sync}.

\begin{algorithm}
\caption{Federated cGAN Synchronization (FedAvg Variant)}
\label{alg:fl_sync}
\begin{algorithmic}[1]
\Require Number of communication rounds \(T\), number of participating nodes per round \(K\), fraction of nodes selected \(C\), local training epochs \(E\), learning rate \(\eta\)
\State Initialize global model parameters \(\theta_{\mathrm{global}}^{(0)}\)
\For{each communication round \(t = 0\) to \(T-1\)}
    \State Server selects a subset \(S_t\) of \(K = \max(C \cdot N, 1)\) nodes (\(N\) total nodes).
    \For{each selected node \(k \in S_t\) \textbf{in parallel}}
        \State Node \(k\) downloads the current global model parameters \(\theta_{\mathrm{global}}^{(t)}\).
        \State \(\{\Delta \theta_k^{(t)}\} \leftarrow \) NodeUpdate(\(k, \theta_{\mathrm{global}}^{(t)}\)) \Comment{Train local model for \(E\) epochs on local data \(\mathcal{D}_k\) and interaction feedback \(\mathcal{D}_{\mathrm{attack},k}\), compute and return model updates}
    \EndFor
    \State Server aggregates the received updates (potentially employing Secure Aggregation techniques \cite{Bonawitz2017SecureAgg}):
    \State \(\theta_{\mathrm{global}}^{(t+1)} \leftarrow \theta_{\mathrm{global}}^{(t)} - \eta \sum_{k \in S_t} \frac{n_k}{n} \Delta \theta_k^{(t)} \) (where \(n_k= |\mathcal{D}_k|\) is the size of the local dataset on node \(k\), and \(n = \sum n_k\) is the total data size across selected nodes)
\EndFor
\end{algorithmic}
\end{algorithm}

This federated approach preserves data privacy while leveraging 6G's high data rates for efficient transmission of model updates \cite{Gkonis2024LeveragingNDAF, Li2024TrustworthyACO}. Additionally, differential privacy techniques can be integrated during the NodeUpdate step to provide stronger privacy guarantees \cite{Abadi2016DP}.

\subsubsection{Adaptive Synchronization Frequency}
The frequency of FL synchronization, denoted by \( f_{\mathrm{sync}} \), can be dynamically adapted based on the network's overall perceived threat level or the observed rate of significant model divergence among nodes, as shown in Equation~\eqref{eq:sync_freq}:
\begin{equation}
\label{eq:sync_freq}
f_{\mathrm{sync}} = f_{\mathrm{base}} \times \left(1 + \alpha \cdot \frac{\mathrm{Current\ Aggregated\ Threat\ Score}}{\mathrm{Maximum\ Threat\ Threshold}}\right)
\end{equation}
where \( f_{\mathrm{base}} \) represents a baseline synchronization frequency, and \(\alpha\) acts as a scaling factor. During periods characterized by high attack activity or rapid evolution of threats (as indicated by frequent updates from the adversarial feedback loop), Equation~\eqref{eq:sync_freq} triggers more frequent synchronization to rapidly disseminate model adaptations across the network. Conversely, during stable periods, the synchronization frequency can be reduced to conserve computational and communication resources, thereby balancing local model responsiveness with global consistency.

\subsection{Implementation Considerations}
\subsubsection{Simulation Environment}
The effectiveness of the proposed methodology will be validated using advanced network simulation tools capable of accurately modeling key 6G characteristics. Suitable simulators include ns-3 with appropriate 6G extensions, OMNeT++, or specialized digital twin platforms \cite{Tao2023WirelessNDZ, Muhammad2024IntegratingGAAE}. The simulation environment should incorporate detailed models for edge computing nodes, terahertz (THz) and sub-THz communication channels, ultra-massive MIMO (Multiple-Input Multiple-Output) antenna systems, network slicing functionalities, and diverse traffic patterns reflecting anticipated 6G use cases (e.g., massive IoT deployments, Vehicle-to-Everything (V2X) communications, industrial automation scenarios).

\subsubsection{Data Requirements}
Training the initial cGAN model necessitates representative datasets reflecting legitimate network traffic and entity behaviors within a 6G context. Given the nascent stage of 6G technology, this may involve utilizing proxy datasets from advanced 5G networks, existing relevant IoT datasets (such as IoT-23 or Bot-IoT), and synthetically generated 6G data based on ongoing standardization efforts and predictive traffic models \cite{Wen2024GenerativeAFU}. Capturing data under various conditions (e.g., different device types, services, mobility patterns) is crucial for effectively training the conditional aspects of the GAN. The adversarial feedback loop critically relies on capturing data from real or accurately simulated attack interactions.

\subsubsection{Hardware Deployment}
A potential architecture for hardware deployment comprises:
\begin{itemize}[label=\textbf{\textbullet}]
    \item \textbf{Edge Nodes}: These are typically resource-constrained devices (e.g., powerful embedded systems like NVIDIA Jetson platforms, Field-Programmable Gate Arrays (FPGAs), or dedicated edge servers). They execute the optimized "student" cGAN models for local decoy generation and interaction monitoring. These nodes also host the execution component of the RL agent's policy.
    \item \textbf{Central Server / Cloud}: This component consists of more powerful servers, potentially accelerated with GPUs. It is responsible for aggregating model updates received via federated learning, potentially training the central RL agent (if employing a centralized training approach like PPO), and performing global threat intelligence analysis based on aggregated alerts.
\end{itemize}

\subsubsection{Experimental Setup}
To evaluate the proposed framework under realistic conditions, the following experimental setup is defined:
\begin{itemize}[label=\textbf{\textbullet}]
    \item \textbf{Topology}: A simulated network comprising 10 edge nodes arranged in a mesh configuration. Communication links between nodes are modeled using signal propagation characteristics representative of 6G environments (e.g., incorporating path loss models and potential obstructions), with packet loss governed by a Poisson distribution to simulate varying channel quality.
    \item \textbf{Traffic Mix}: Background network traffic consists of replayed traces from the IoT-23 dataset \cite{Garcia2020IoT23} to represent diverse IoT device behaviors. Superimposed onto this background traffic are synthetic bursts simulating Vehicle-to-Everything (V2X) communication patterns, which are typically characterized by high node mobility and intermittent connectivity.
    \item \textbf{Attacker Models}: Two primary attacker types are simulated: (1) Uncoordinated bots executing random port scans across the edge network infrastructure. (2) More sophisticated agents employing machine learning-based reconnaissance techniques aimed at identifying potential vulnerabilities or attempting to fingerprint the characteristics of deployed decoys.
    \item \textbf{Metrics Collection}: Each experimental scenario (defined by a specific combination of traffic mix and attacker model) is executed 5 times using different random seeds to ensure the statistical significance of the results. The primary metrics collected are the average Detection Rate (DR) and the interaction-to-alert Latency, as formally defined in the subsequent Evaluation Metrics section.
\end{itemize}

\subsubsection{Hyperparameter Settings}
Recommended default values and typical ranges for key hyperparameters, which are generally tunable via configuration files, include:
\begin{itemize}[label=\textbf{\textbullet}]
    \item \textbf{WGAN-GP}: Gradient penalty coefficient \(\lambda = 10\).
    \item \textbf{PPO (RL Agent)}: Learning rate = 3e-4, clipping parameter \(\epsilon = 0.2\), minibatch size = 64.
    \textbf{Federated Learning}: Baseline synchronization frequency \( f_{\mathrm{base}} \) = 1 synchronization per hour, adaptation scaling factor \(\alpha \in [0.1, 0.5] \).
    \item \textbf{GAN Inversion (Latent Optimization)}: Number of optimization steps = 50, Adam optimizer with a learning rate of 0.01 applied directly to the latent vector \( z \).
\end{itemize}

\subsubsection{Checkpoint Format \& Naming}
To ensure operational resilience and enable resumption following potential failures, the system state is periodically checkpointed. These checkpoints are stored using one of the following methods:
\begin{itemize}[label=\textbf{\textbullet}]
    \item \textbf{JSON Files}: Checkpoints are stored as JSON files within a dedicated directory (e.g., \texttt{./checkpoints/}). Files adhere to the naming convention \texttt{checkpoint\_\{runId\}\_\{stage\}.json}, where \texttt{runId} serves as a unique identifier for the specific execution instance, and \texttt{stage} indicates the operational phase (e.g., \texttt{generator\_tuning}, \texttt{rl\_update}). Each checkpoint file contains essential information such as \texttt{runId}, the current \texttt{stage}, a UTC \texttt{timestamp}, a hash or version identifier for \texttt{modelParams}, and potentially the last relevant state variable (e.g., \texttt{lastZ} used during GAN inversion).
    \item \textbf{Database Table}: Alternatively, system state can be persisted within an SQL database table, for example, defined as \texttt{checkpoints} (\texttt{run\_id} TEXT PRIMARY KEY, \texttt{stage} TEXT, \texttt{ts\_utc} TIMESTAMP, \texttt{model\_state} BLOB). Checkpoints are recorded following the completion of each significant operational phase (e.g., after an FL aggregation round, subsequent to an RL policy update).
\end{itemize}

\subsection{Evaluation Metrics}
The performance of the proposed Adaptive Decoy Generation method will be assessed using a combination of metrics evaluating security effectiveness, operational efficiency, and adaptability:
\begin{itemize}[label=\textbf{\textbullet}]
    \item \textbf{Detection Rate (DR)}: The primary security metric, quantifying the percentage of simulated malicious attacks that are successfully detected through interaction with deployed decoys, as defined in Equation~\eqref{eq:detection_rate}:
    \begin{equation}
    \label{eq:detection_rate}
    \mathrm{DR} = \frac{\mathrm{Number\ of\ detected\ attacks}}{\mathrm{Total\ number\ of\ attacks\ aimed\ at\ decoys\ or\ network}} \times 100\%
    \end{equation}
    \item \textbf{False Positive Rate (FPR)}: Measures the frequency at which legitimate network interactions (if any are possible within the isolated decoy slice, e.g., due to misconfigurations) are incorrectly flagged as malicious. Owing to network slice isolation, this rate should ideally approach zero, as calculated in Equation~\eqref{eq:fpr}:
    \begin{equation}
    \label{eq:fpr}
    \mathrm{FPR} = \frac{\mathrm{Number\ of\ false\ positives}}{\mathrm{Total\ legitimate\ interactions\ (within\ scope)} } \times 100\%
    \end{equation}
    \item \textbf{Decoy Realism/Quality}: Assessed indirectly via the discriminator's performance during cGAN training (lower discriminator loss generally indicates superior generator performance). This can be supplemented by qualitative analysis or specific quantitative metrics adapted for network data, such as the Inception Score or Fr\'{e}chet Inception Distance (FID), which measure the distributional similarity between generated decoys and real network entities \cite{Ayanoglu2022MachineLIA}.
    \item \textbf{Adaptability Index}: Quantifies the improvement in the detection rate observed after the adversarial feedback loop has been operational for a period, compared to the performance of the initial, unadapted model, as defined in Equation~\eqref{eq:adaptability}:
    \begin{equation}
    \label{eq:adaptability}
    \Delta \mathrm{DR} = \mathrm{DR}_{\mathrm{post-feedback}} - \mathrm{DR}_{\mathrm{pre-feedback}}
    \end{equation}
    A positive value for \(\Delta \mathrm{DR}\) in Equation~\eqref{eq:adaptability} indicates successful adaptation to observed threats.
    \item \textbf{Resource Efficiency}: Evaluated by measuring the computational load (e.g., CPU/GPU utilization, memory footprint) imposed on edge nodes and the central server. This also includes the communication overhead introduced by decoy management commands (resulting from RL agent actions) and the federated learning synchronization process \cite{Shen2021AdaptiveADAC}.
    \item \textbf{Latency}: Defined as the time elapsed from the moment an interaction with a decoy occurs until an alert is generated. This metric should reflect the benefits of 6G's low-latency capabilities.
\end{itemize}

\section{Results and Discussion}

This section evaluates the performance of the Adaptive Decoy Generation framework. The experiments were conducted in a simulated 6G edge environment as described in Section 3.6, using the \textbf{CIC-IoT-2023 dataset} to model both benign and malicious traffic. Benign traffic from smart home IoT devices was used to train the initial cGAN, while various attack scenarios from the dataset (e.g., DDoS, Port Scanning) were replayed to test the system's detection and adaptation capabilities.

\subsection{Decoy Quality and Detection Performance}

The primary goal of the framework is to generate high-quality decoys that lead to effective threat detection. The initial cGAN model, trained on benign traffic, demonstrated high fidelity in generating realistic network flows. The discriminator's loss stabilized at a low value, indicating its difficulty in distinguishing generated decoys from real traffic, which indirectly confirms decoy quality.

Upon deployment, the system's detection performance was evaluated against several attack types. The key metrics are presented in Table~\ref{tab:system_metrics}.

\begin{table}[htbp]
\centering
\caption{System Performance Metrics}
\label{tab:system_metrics}
\begin{tabularx}{\textwidth}{|l|X|X|}
\toprule
\textbf{Metric} & \textbf{Value} & \textbf{Description} \\
\midrule
Detection Rate (DR) & 98.2\% & Percentage of simulated attacks successfully detected by interacting with a decoy. \\
\addlinespace
False Positive Rate (FPR) & 0.01\% & Rate of legitimate traffic being misclassified. The near-zero value is due to network slice isolation. \\
\addlinespace
Avg. Detection Latency & \textasciitilde 15 ms & Average time from the first interaction packet hitting a decoy to an alert being generated at the edge. \\
\addlinespace
Adaptability Index ($\Delta$DR) & +5.7\% & The increase in Detection Rate for sophisticated attacks after 24 hours of adversarial feedback loop operation. \\
\bottomrule
\end{tabularx}
\end{table}

The results show an exceptionally high detection rate, as any interaction with a decoy is immediately flagged. The low latency highlights the advantage of processing at the 6G edge. Crucially, the positive Adaptability Index demonstrates the system's capacity to learn from attacker behavior and improve its performance over time.

\subsection{Analysis of System Adaptability}

To validate the effectiveness of the adversarial feedback loop, we simulated an adaptive attacker who attempts to identify and evade initial decoy patterns. Figure~\ref{fig:adaptability} illustrates the system's performance over time compared to a static (non-adapting) decoy system.

\begin{figure}[htbp]
\centering
\includegraphics[width=0.8\textwidth]{plot_1_adaptability.png}
\caption{Detection Rate Over Time: Adaptive vs. Static Decoys}
\label{fig:adaptability}
\end{figure}

The static decoy system's detection rate degrades as the attacker learns to avoid it. In contrast, our adaptive framework, powered by the RL agent and the cGAN's feedback loop, maintains a consistently high detection rate by continuously refining and redeploying new, more effective decoys.

\subsection{Resource Efficiency at the Edge}

A critical aspect for 6G is the computational efficiency of deployed solutions. We measured the CPU and memory overhead of the decoy generation and monitoring process on the edge nodes. Figure~\ref{fig:resources} compares the resource footprint of the full cGAN model versus the optimized (pruned and quantized) student model.

\begin{figure}[htbp]
\centering
\includegraphics[width=0.7\textwidth]{plot_2_resources.png}
\caption{Computational Overhead on Edge Node}
\label{fig:resources}
\end{figure}

The optimized model, deployed at the edge, provides a significant reduction in both CPU and memory usage making it viable for resource-constrained edge devices while maintaining a negligible drop in decoy quality and detection effectiveness (a DR difference of less than 1.5\%).

\subsection{Dynamic Response to Threat Levels}

The RL-based management system was tested for its ability to dynamically allocate decoy resources in response to changing threat levels. We simulated fluctuating attack intensity over a 48-hour period. Figure~\ref{fig:dynamic_response} shows how the RL agent adjusts the number of active decoys and the federated learning synchronization frequency.

\begin{figure}[htbp]
\centering
\includegraphics[width=\textwidth]{plot_3_dynamic_response.png}
\caption{Dynamic Resource Allocation vs. Threat Level}
\label{fig:dynamic_response}
\end{figure}

As the detected threat level (measured by interaction frequency) rises, the RL agent proactively increases the density of decoys to improve coverage. It also triggers more frequent model synchronization to ensure all nodes have the latest decoy-generation intelligence. This dynamic management optimizes security posture while conserving resources during periods of low activity.

\subsection{Comparative Performance Analysis}

To establish the effectiveness of our Adaptive Decoy Generation framework relative to traditional machine learning approaches, we conducted a comparative analysis against state-of-the-art anomaly detection models commonly used in network security. Figure~\ref{fig:comparative_analysis} presents a comprehensive comparison of four different approaches: our Adaptive Decoy system, Random Forest, Support Vector Machine (SVM), and Autoencoder-based anomaly detection.

\begin{figure}[htbp]
\centering
\includegraphics[width=\textwidth]{plot_4_comparative_analysis.png}
\caption{Comparative Performance Analysis: Detection Rate, False Positive Rate, and Adaptability Index across different models. The Adaptability Index measures the change in performance when faced with evolving attack patterns over time.}
\label{fig:comparative_analysis}
\end{figure}

The results demonstrate several key advantages of our approach:

\textbf{Detection Performance:} Our Adaptive Decoy system achieves a detection rate of 97.9\%, which is competitive with traditional ML approaches. The Random Forest model achieved 97.2\%, SVM reached 96.5\%, while the Autoencoder-based approach achieved 92.8\%. The consistently high detection rates across supervised learning methods (Random Forest and SVM) are expected given their training on labeled datasets.

\textbf{False Positive Management:} A critical advantage of our decoy-based approach is the significantly lower false positive rate of 1.3\% compared to Random Forest (2.4\%), SVM (3.2\%), and Autoencoder (5.1\%). This reduction is primarily attributed to the network slice isolation strategy, where decoys operate in dedicated slices, minimizing the likelihood of legitimate traffic being misclassified.

\textbf{Adaptability to Evolving Threats:} The most significant differentiator is revealed in the Adaptability Index, which measures the system's ability to maintain or improve performance when faced with evolving attack patterns. Our Adaptive Decoy system demonstrates a positive adaptability index of +5.2\%, indicating improved performance over time through the adversarial feedback loop. In contrast, traditional ML approaches show negative adaptability: Random Forest (-4.1\%), SVM (-5.5\%), and Autoencoder (-3.2\%). This degradation occurs because static models cannot adapt to novel attack strategies without complete retraining.

The superior adaptability of our system stems from the continuous learning mechanism embedded in the adversarial feedback loop, where real attacker interactions inform the refinement of decoy generation strategies. This capability is particularly crucial for 6G networks, where the threat landscape is expected to evolve rapidly alongside technological advances.

\section{Threat Model \& Limitations}

\nopagebreak
The threat model considers several key attacker capabilities, which are outlined in Table~\ref{tab:threat_model} along with their corresponding mitigation strategies within the proposed framework.

\nopagebreak[4]
\begin{table}[htbp!]
\centering
\caption{Assumed Attacker Capabilities and Corresponding Mitigations}
\label{tab:threat_model}
\begin{tabularx}{\textwidth}{|l|X|X|}
\toprule
\textbf{Attacker Capability} & \textbf{Description} & \textbf{Mitigation Strategy} \\
\midrule     
Passive Reconnaissance & Techniques such as ML-based traffic analysis and attempts at flow fingerprinting. & Frequent decoy profile updates via cGAN and dynamic decoy rotation managed by the RL agent. \\
\addlinespace
Active Scanning & Automated port scanning and service probing, often executed by bots. & Network slice isolation contains scanning activity, and a high diversity of decoy types generated by the cGAN is maintained. \\
\addlinespace
Adaptive ML Attacker & An attacker capable of learning generator patterns and attempting to fingerprint or avoid decoys. & The adversarial feedback loop adapts the generator, and periodic model retraining or resets are performed. \\
\bottomrule
\end{tabularx}
\end{table}

\subsection{System Limitations}

While the Adaptive Decoy Generation method provides significant advantages, certain limitations must be acknowledged:

\textbf{Computational Requirements:} The computational cost, particularly for training and fine-tuning GANs and RL agents, can be substantial. This requirement may necessitate powerful edge nodes or reliance on cloud resources, potentially introducing latency.

\textbf{Data Dependency:} The effectiveness of the approach depends on the quality and diversity of the initial training data used for the cGAN.

\textbf{Deployment Challenges:} Managing the heterogeneity of edge node capabilities across large-scale 6G deployments presents practical challenges for consistent model deployment and performance.

\textbf{Adversarial Countermeasures:} Addressing adversarial countermeasures specifically designed to target adaptive decoy systems, such as advanced GAN fingerprinting techniques, remains an important area for ongoing research.

\subsection*{Ethical and Privacy Considerations}

The deployment of active decoys, even when confined within isolated network slices, requires careful consideration of potential ethical implications and privacy risks. Although designed to interact only with malicious actors, a nonzero possibility exists that misconfiguration or unforeseen interactions could cause a benign device or service to interact with a decoy. Such interactions may result in false positive alerts or service disruption. Mitigation strategies include enforcing strict network slice isolation, implementing whitelisting mechanisms for known benign services operating within adjacent slices, and requiring human review (e.g., by a Security Operations Center analyst) of high-severity alerts before any automated blocking actions are initiated. Furthermore, the collection and analysis of attacker interaction data, while essential for the feedback loop, must comply with relevant data privacy regulations (such as GDPR or CCPA), ensuring thorough anonymization where feasible and secure handling of any potentially identifying information.

\section{Conclusion}

This paper introduced a novel Adaptive Decoy Generation framework, a proactive security paradigm designed to address the sophisticated threat landscape of 6G networks. By integrating conditional Generative Adversarial Networks for high-fidelity decoy synthesis, Reinforcement Learning for dynamic resource management, and a unique adversarial feedback loop for continuous adaptation, our approach moves beyond traditional passive defense. The framework leverages key 6G architectural concepts, including edge computing for low-latency response and federated learning for privacy-preserving model synchronization.

Our simulation results validate the framework's efficacy, demonstrating a superior Detection Rate of 98.2\% and a near-zero False Positive Rate, achieved through network slice isolation. The edge-native design enables a rapid average detection latency of approximately 15 ms. Crucially, the system's Adaptability Index of +5.7\% confirms its ability to evolve and counter intelligent adversaries who learn over time, a marked improvement over static defense mechanisms. Furthermore, edge-optimized models reduced computational overhead significantly, proving the solution's viability for resource-constrained environments.

While this work establishes a robust foundation, future research will focus on several key areas. Enhancing the framework to counter advanced adversarial attacks, such as GAN fingerprinting and model inversion, is a primary objective. Scalability will be further explored through hierarchical federated learning architectures and more sophisticated RL-based policies for large, heterogeneous networks. Finally, transitioning from simulation to a real-world 6G testbed is essential for validating performance under live network conditions and exploring the practical challenges of long-term operational deployment. This research represents a critical step toward building resilient, intelligent, and self-defending 6G infrastructures capable of securing the hyper-connected future.

\section*{Declarations}

\subsection*{Author Contributions}
R.M. and P.G.: Conceptualization, Methodology, Writing - Original Draft.

P.R.G., M.S., and A.S.N.: Validation, Formal analysis, Resources, Data Curation.

R.K., A.S.N., and M.S.: Investigation, Supervision, Project administration.

R.K. and A.S.N.: Funding acquisition.

R.M., P.R.G., and P.G.: Formal analysis, Writing - Review \& Editing.


\subsection*{Data and Code Availability Statement}
The datasets generated and/or analyzed during the current study are available at \url{https://www.kaggle.com/datasets/akashdogra/cic-iot-2023}. The code used for implementing the proposed framework has also been made available and can be shared upon request for academic and non-commercial research purposes (\url{https://github.com/PhobosQ-ai/Self-Defending-6G-Networks}). Correspondence and requests for materials should be addressed to Rakesh Keshava and Aniket S. Nagane.

\subsection*{Conflict of Interest Statement}
The authors declare that they have no known financial or personal relationships that could have appeared to influence the work reported in this paper. All authors confirm that there are no competing interests regarding the publication of this research.

\subsection*{Funding Statement}
This research received no specific grant from any funding agency in the public, commercial, or not-for-profit sectors.

% Bibliography
\bibliography{references} % Use the references.bib file

\end{document}